\documentclass[10pt,a4paper]{amsart}
\usepackage{amsmath,amssymb,amsthm,graphicx}
% Retained minimal preamble for the Stacks Project chapter varieties.tex.
% Verbatim excerpt of the chapter preamble: Xy-pic setup (native preamble lines 14-18), the native theorem
% declarations with the mathematical macros (native preamble lines 154-217), and the project own no-op
% definitions of the reference/slogan/history display environments (native wrapper.tex lines 1-3) that the
% counted statement uses for its slogan. Copied mechanically; nothing added, removed or redefined.
% For commutative diagrams we use Xy-pic
\usepackage[all]{xy}

% We use 2cell for 2-commutative diagrams.
\xyoption{2cell}
% Theorem environments.
%
\theoremstyle{plain}
\newtheorem{theorem}[subsection]{Theorem}
\newtheorem{proposition}[subsection]{Proposition}
\newtheorem{lemma}[subsection]{Lemma}

\theoremstyle{definition}
\newtheorem{definition}[subsection]{Definition}
\newtheorem{example}[subsection]{Example}
\newtheorem{exercise}[subsection]{Exercise}
\newtheorem{situation}[subsection]{Situation}

\theoremstyle{remark}
\newtheorem{remark}[subsection]{Remark}
\newtheorem{remarks}[subsection]{Remarks}

\numberwithin{equation}{subsection}

% Macros
%
\def\lim{\mathop{\mathrm{lim}}\nolimits}
\def\colim{\mathop{\mathrm{colim}}\nolimits}
\def\Spec{\mathop{\mathrm{Spec}}}
\def\Hom{\mathop{\mathrm{Hom}}\nolimits}
\def\Ext{\mathop{\mathrm{Ext}}\nolimits}
\def\SheafHom{\mathop{\mathcal{H}\!\mathit{om}}\nolimits}
\def\SheafExt{\mathop{\mathcal{E}\!\mathit{xt}}\nolimits}
\def\Sch{\mathit{Sch}}
\def\Mor{\mathop{\mathrm{Mor}}\nolimits}
\def\Ob{\mathop{\mathrm{Ob}}\nolimits}
\def\Sh{\mathop{\mathit{Sh}}\nolimits}
\def\NL{\mathop{N\!L}\nolimits}
\def\CH{\mathop{\mathrm{CH}}\nolimits}
\def\proetale{{pro\text{-}\acute{e}tale}}
\def\etale{{\acute{e}tale}}
\def\QCoh{\mathit{QCoh}}
\def\Ker{\mathop{\mathrm{Ker}}}
\def\Im{\mathop{\mathrm{Im}}}
\def\Coker{\mathop{\mathrm{Coker}}}
\def\Coim{\mathop{\mathrm{Coim}}}

% Boxtimes
%
\DeclareMathSymbol{\boxtimes}{\mathbin}{AMSa}{"02}

%
% Macros for moduli stacks/spaces
%
\def\QCohstack{\mathcal{QC}\!\mathit{oh}}
\def\Cohstack{\mathcal{C}\!\mathit{oh}}
\def\Spacesstack{\mathcal{S}\!\mathit{paces}}
\def\Quotfunctor{\mathrm{Quot}}
\def\Hilbfunctor{\mathrm{Hilb}}
\def\Curvesstack{\mathcal{C}\!\mathit{urves}}
\def\Polarizedstack{\mathcal{P}\!\mathit{olarized}}
\def\Complexesstack{\mathcal{C}\!\mathit{omplexes}}
% \Pic is the operator that assigns to X its picard group, usage \Pic(X)
% \Picardstack_{X/B} denotes the Picard stack of X over B
% \Picardfunctor_{X/B} denotes the Picard functor of X over B
\def\Pic{\mathop{\mathrm{Pic}}\nolimits}
\def\Picardstack{\mathcal{P}\!\mathit{ic}}
\def\Picardfunctor{\mathrm{Pic}}
\def\Deformationcategory{\mathcal{D}\!\mathit{ef}}
\newenvironment{reference}{}{}
\newenvironment{slogan}{}{}
\newenvironment{history}{}{}
\usepackage{cleveref}
\crefname{originallemma}{original lemma}{original lemmas}
\Crefname{originallemma}{Original lemma}{Original lemmas}
\makeatletter
\newlabel{algebra-lemma-field-extension-geometrically-irreducible}{{037P}{source}{}{sourcepointer-1}{}}
\newlabel{algebra-lemma-field-extension-geometrically-irreducible@cref}{{[originallemma][0][]037P}{[1][0][]source}{}{}{}}
\makeatother
\crefname{originallemma}{original lemma}{original lemmas}
\Crefname{originallemma}{Original lemma}{Original lemmas}
\makeatletter
\newlabel{lemma-characterize-geometrically-connected}{{0387}{source}{}{sourcepointer-2}{}}
\newlabel{lemma-characterize-geometrically-connected@cref}{{[originallemma][0][]0387}{[1][0][]source}{}{}{}}
\makeatother
\crefname{originallemma}{original lemma}{original lemmas}
\Crefname{originallemma}{Original lemma}{Original lemmas}
\makeatletter
\newlabel{lemma-geometrically-connected-criterion}{{056R}{source}{}{sourcepointer-3}{}}
\newlabel{lemma-geometrically-connected-criterion@cref}{{[originallemma][0][]056R}{[1][0][]source}{}{}{}}
\makeatother
\title{Connectedness plus a point with algebraically closed residue field gives geometric connectedness}
\author{The Stacks Project Authors}
\date{stacks-project:varieties.tex:local-snapshot-20261011; GNU FDL 1.2 or later}
\begin{document}
\maketitle
\noindent\textit{Editorial note.} All delivered mathematical text is the author's own, copied from the pinned Stacks Project chapter file varieties.tex; no statement or proof step was rewritten, reordered or removed.\par\medskip
\noindent\textit{Editorial note.} The proof is complete as written and ends with the characterization applied to X; no step is referred to the reader.\par\medskip
\noindent\textit{Editorial note.} The counted spans carry 3 navigation stubs and no bibliography citation.\par\medskip
\noindent\textit{Editorial note.} The preamble is the chapter's retained minimal preamble (sources/stacks-varieties/retained-preamble.tex lines 1-77), a verbatim excerpt of the native preamble.\par\medskip
\noindent\textit{Editorial note.} The label lemma-geometrically-connected-if-connected-and-point is defined only in varieties.tex, so no cross-chapter restatement had to be weighed for this card.\par\medskip
\noindent\textit{Allowed standard prerequisites (editorial).}\begin{itemize}
\item A field extension in which k is algebraically closed becomes geometrically irreducible after base change. Algebra, Lemma algebra-lemma-field-extension-geometrically-irreducible, invoked at proof lines 1093-1095. Supplies irreducibility of the base change of the spectrum of the residue field.
\item Criterion for geometric connectedness through a connected morphism to an irreducible base. Same chapter, Lemma lemma-geometrically-connected-criterion, invoked at proof lines 1096-1098. Supplies connectedness of X over the separable closure.
\item Characterization of geometric connectedness by connectedness over the separable closure. Same chapter, Lemma lemma-characterize-geometrically-connected, invoked at proof lines 1098-1100. Supplies the final conclusion.
\end{itemize}\medskip
\ifdefined\hypertarget\hypertarget{sourcepointer-1}{}\fi
\noindent\textit{Source reference.} Original lemma 037P: sources/stacks-chow/tags.txt, source lines 4208--4208. Pointer outside the counted span, invoked once for context; the load-bearing argument of the counted statement is the author's proof reproduced in full in the proof span.\par\medskip
\ifdefined\hypertarget\hypertarget{sourcepointer-2}{}\fi
\noindent\textit{Source reference.} Original lemma 0387: sources/stacks-varieties/pinned-original-source.tex, source lines 792--792. Pointer outside the counted span, invoked once for context; the load-bearing argument of the counted statement is the author's proof reproduced in full in the proof span.\par\medskip
\ifdefined\hypertarget\hypertarget{sourcepointer-3}{}\fi
\noindent\textit{Source reference.} Original lemma 056R: sources/stacks-varieties/pinned-original-source.tex, source lines 1067--1067. Pointer outside the counted span, invoked once for context; the load-bearing argument of the counted statement is the author's proof reproduced in full in the proof span.\par\medskip
\begin{lemma}
\label{lemma-geometrically-connected-if-connected-and-point}
Let $k$ be a field. Let $X$ be a scheme over $k$.
Assume $X$ is connected and has a point $x$ such that
$k$ is algebraically closed in $\kappa(x)$.
Then $X$ is geometrically connected.
In particular, if $X$ has a $k$-rational point and $X$ is connected,
then $X$ is geometrically connected.
\end{lemma}

\begin{proof}
Set $T = \Spec(\kappa(x))$. Let $\overline{k}$ be a
separable algebraic closure of $k$. The assumption on $\kappa(x)/k$
implies that $T_{\overline{k}}$ is irreducible, see
Algebra, Lemma \ref{algebra-lemma-field-extension-geometrically-irreducible}.
Hence by
Lemma \ref{lemma-geometrically-connected-criterion}
we see that $X_{\overline{k}}$ is connected. By
Lemma \ref{lemma-characterize-geometrically-connected}
we conclude that $X$ is geometrically connected.
\end{proof}

\end{document}
